\heading{数学物理方法（下）-第8次作业}


\begin{question}
利用本征值问题证明连带 Legendre 函数的正交性
\[
\int_{-1}^1P_l^m(x)P_k^m(x)\,dx=0,\qquad k\ne l.
\]
\begin{solution}
\(P_l^m\) 满足 Sturm--Liouville 方程
\[
\frac d{dx}\left[(1-x^2)y'\right]
+\left[l(l+1)-\frac{m^2}{1-x^2}\right]y=0.
\]
对 \(l,k\) 两个方程分别乘以 \(P_k^m,P_l^m\) 后相减并积分，端点项因有界性消失，得
\[
\{l(l+1)-k(k+1)\}\int_{-1}^{1}P_l^mP_k^m\,dx=0.
\]
当 \(l\ne k\) 时即得正交性。
\end{solution}
\end{question}

\begin{question}
一半径为 \(a\) 的均匀导体球，表面温度为
\[
u|_{r=a}=P_1^1(\cos\theta)\cos\phi,
\]
球内无热源，求球内稳态温度分布。
\begin{solution}
球内有界调和函数为
\[
u=\sum_{l,m}r^l\left(A_{lm}\cos m\phi+B_{lm}\sin m\phi\right)P_l^m(\cos\theta).
\]
边界只含 \(l=1,m=1\) 的余弦项，因此
\[
u(r,\theta,\phi)=\frac ra P_1^1(\cos\theta)\cos\phi.
\]
\end{solution}
\end{question}

\begin{question}
将 \(\sin^2\theta\cos^2\phi\) 与 \((1+\cos\theta)\sin\theta\cos\phi\) 按球面调和函数展开。
\begin{solution}
使用实球谐函数更简洁。由
\[
\sin^2\theta\cos^2\phi=\frac13-\frac13P_2(\cos\theta)
+\frac13P_2^2(\cos\theta)\cos2\phi,
\]
以及
\[
(1+\cos\theta)\sin\theta\cos\phi
=-P_1^1(\cos\theta)\cos\phi-\frac13P_2^1(\cos\theta)\cos\phi,
\]
即可换成标准归一化的 \(Y_l^m\) 展开；只需按所采用的归一化常数乘除。
\end{solution}
\end{question}

\begin{question}
求解球内问题
\[
\nabla^2u=0,\qquad r<a,\qquad
u|_{r=a}=A+B\sin2\theta\cos\phi.
\]
\begin{solution}
因
\[
\sin2\theta\cos\phi=2\sin\theta\cos\theta\cos\phi
=-\frac23P_2^1(\cos\theta)\cos\phi,
\]
球内有界解为
\[
u(r,\theta,\phi)=A-\frac{2B}{3}\left(\frac ra\right)^2
P_2^1(\cos\theta)\cos\phi.
\]
\end{solution}
\end{question}

\begin{question}
求解球内 Poisson 问题
\[
\nabla^2u=A+Br^2\sin2\theta\cos\phi,\qquad r<a,\qquad u|_{r=a}=0.
\]
\begin{solution}
常数源项给出径向特解
\[
u_0=\frac A6(r^2-a^2).
\]
又
\[
r^2\sin2\theta\cos\phi=-\frac23r^2P_2^1(\cos\theta)\cos\phi.
\]
利用
\[
\nabla^2\left(r^{l+2}Y_l^m\right)=2(2l+3)r^lY_l^m,
\]
取 \(l=2\)，得源项特解
\[
u_1=-\frac{B}{21}r^4P_2^1(\cos\theta)\cos\phi.
\]
为使边界为零，补上同阶调和项：
\[
u(r,\theta,\phi)=\frac A6(r^2-a^2)
\frac{B}{21}(a^2r^2-r^4)P_2^1(\cos\theta)\cos\phi.
\]
\end{solution}
\end{question}

